\SourceLeaf{446}%
\begin{quote}
»Zieht man von einem Punkt der krummen Fläche aus lauter kürzeste Linien gleicher Länge, so bilden diese mit der Linie, die ihre Endpunkte verbindet, überall rechte Winkel«.
\end{quote}
\noindent\hspace*{1.2em}Ganz ähnlich lässt sich Folgendes beweisen. Auf der krummen Fläche sei eine beliebige krumme Linie gegeben. Gehen von jedem ihrer Punkte nach derselben Seite und im rechten Winkel zu ihr kürzeste Linien gleicher Länge aus, so schneiden diese die Linie, die ihre Endpunkte verbindet, überall im rechten Winkel. Man braucht in der vorigen Herleitung $\theta$ nur als die Länge auf der gegebenen krummen Linie von einem beliebigen Punkt aus zu verstehen. Die vorigen Rechnungen bleiben dann gültig, wobei sich der Wert $u=0$ für $s=0$ nun schon aus der Voraussetzung ergibt.

\Needspace{5\baselineskip}
\begin{center}18.\end{center}
\noindent\hspace*{1.2em}Die Größenverhältnisse bei diesen Konstruktionen verdienen es, ausführlicher hergeleitet zu werden. Schreibt man der Kürze halber $m$ für $\frac{\partial t}{\partial\theta}$, so gilt zunächst:
\begin{equation}
\frac{\partial x}{\partial s}=\xi,\qquad\frac{\partial y}{\partial s}=\eta,\qquad\frac{\partial z}{\partial s}=\zeta,\tag{1}
\end{equation}
\begin{equation}
\frac{\partial x}{\partial\theta}=m\xi',\qquad\frac{\partial y}{\partial\theta}=m\eta',\qquad\frac{\partial z}{\partial\theta}=m\zeta',\tag{2}
\end{equation}
\begin{equation}\xi\xi+\eta\eta+\zeta\zeta=1\tag{3}\end{equation}
\begin{equation}\xi'\xi'+\eta'\eta'+\zeta'\zeta'=1\tag{4}\end{equation}
\begin{equation}\xi\xi'+\eta\eta'+\zeta\zeta'=0.\tag{5}\end{equation}
Ferner gilt
\begin{equation}XX+YY+ZZ=1\tag{6}\end{equation}
\begin{equation}X\xi+Y\eta+Z\zeta=0\tag{7}\end{equation}
\begin{equation}X\xi'+Y\eta'+Z\zeta'=0\tag{8}\end{equation}
und
\begin{equation}
\left\{\begin{aligned}X&=\zeta\eta'-\eta\zeta'\\Y&=\xi\zeta'-\zeta\xi'\\Z&=\eta\xi'-\xi\eta'\end{aligned}\right.\tag*{[9]}
\end{equation}
